%!TEX root = ../main.tex
\chapter{Introduction}
\label{ch:intro}

Automatic modulation classification (AMC) is a fundamental capability in
software-defined radio, cognitive radio networks, and spectrum monitoring
systems~\cite{dobre2007survey,hazza2013overview}. By identifying the
modulation format of received signals without prior coordination with the
transmitter, AMC enables dynamic spectrum access, interference management,
and electronic surveillance. Recent deep-learning approaches---including
convolutional networks~\cite{otoole2016rml}, CLDNN
architectures~\cite{zhang2023cldnn}, and the Adaptive Wavelet Network
(AWN)~\cite{li2023awn}---achieve classification accuracy above 91\% on the
standard RML2016.10a benchmark across 11 modulation classes. However, the same
properties that make deep neural network (DNN) classifiers effective render
them vulnerable to adversarial perturbations: small, carefully crafted
modifications to in-phase/quadrature (IQ) signals can collapse classification
accuracy to near
0\%~\cite{szegedy2014intriguing,silvaco2023adversarial_rf,lin2020tactics}.

The adversarial threat to AMC systems is qualitatively different from the
image-classification setting that spawned most adversarial machine-learning
research~\cite{goodfellow2015fgsm,madry2018pgd}. RF IQ signals have typical
amplitudes of approximately $\pm 0.02$, so even small $L_\infty$ perturbations
can constitute a large fraction of the signal dynamic range. Optimization-based
attacks such as the Carlini--Wagner ($L_2$) attack
(CW)~\cite{carlini2017towards} and the Elastic-Net Attack on Deep Neural
Networks (EAD)~\cite{chen2020ead} are particularly effective. At high SNR and
under stronger attack normalization settings, such attacks can drive AMC
accuracy close to chance. In our main evaluation, averaged across SNR
$\geq 0$\,dB, CW reduces AWN accuracy to 35.4\%, while EAD-L1 and EAD-EN
reduce it to 6.6\% and 3.4\% respectively. Existing defenses for AMC fall
into two categories: \emph{adversarial training}, which requires retraining
the classifier on attack-generated examples and must be repeated for each new
attack~\cite{madry2018towards,tramer2017ensemble}; and
\emph{input-transformation defenses}, which preprocess signals before
classification and require no model modification~\cite{guo2018input,xu2018feature}.
To the best of our knowledge, no prior work has performed a systematic
side-by-side comparison between frequency-domain adaptive methods and the
classical signal-processing filters---Kalman, Wiener, Savitzky-Golay, Gaussian,
and FIR---that communications engineers would naturally apply as
baselines~\cite{haykin2002adaptive,proakis2006digital}.

We propose an adaptive spectral defense pipeline that operates entirely in the
frequency domain. \textbf{The core hypothesis of this thesis is that, under
the evaluated attacks, much of the adversarial energy appears in non-dominant
or unstable spectral components relative to the modulation-carrying spectral
support of the signal.} A fast Fourier transform (FFT) followed by retention
of the $K$ dominant magnitude bins therefore tends to remove the
attack-induced spectral artifacts while preserving the modulation signature.
We augment this Top-$K$ operator with a spectral-flatness routing criterion
that detects wideband modulations (e.g., AM-SSB) for which Top-$K$ filtering
would destroy signal content, and an adaptive-$K$ mechanism that automatically
selects per-sample filtering intensity based on the spectral magnitude
knee-point. We benchmark all proposed methods against five per-SNR-calibrated
classical signal-processing filters and an empirical randomized-smoothing
baseline~\cite{cohen2019certified} (majority vote over $k = 20$
Gaussian-noised copies, without computing the formal certified radius),
providing a systematic evaluation of this design space for AMC defense.

This thesis makes four contributions:

\begin{enumerate}
\item \textbf{Control-plane attack characterization.}
  We experimentally characterize adversarial attacks on AMC as
  \emph{control-plane attacks}: under a two-track CRC experiment on the
  strongest optimization-based attack (CW), adversarial perturbations
  cause the classifier to select the wrong demodulator while the
  underlying waveform remains largely decodable under oracle
  demodulation. Across eight digital modulations, CRC integrity under
  oracle demodulation remains above 71\% at SNR 18\,dB, demonstrating
  that the adversarial perturbation leaves the data waveform
  substantially intact. This motivates frequency-domain defense: if the
  perturbation is a spectral artifact, filtering can remove it.
  Generalizing this analysis to EAD, FGSM, and PGD is left as future work.

\item \textbf{Adaptive-$K$: unified adaptive spectral defense.}
  We propose a single receiver-side preprocessing method that requires no
  classifier retraining. Adaptive-$K$ combines spectral-flatness routing
  for wideband modulations, an SNR-adaptive K-cap that adjusts filtering
  aggressiveness to estimated signal quality, and per-sample magnitude-knee
  estimation, all in a single shared-FFT pass.

\item \textbf{Systematic classical-filter comparison.}
  We provide a systematic evaluation of five classical signal-processing
  filters---Kalman, Wiener, Savitzky-Golay, Gaussian, and FIR---as adversarial
  defenses for AMC. Filter parameters are calibrated per-SNR via grid search
  on validation accuracy. The results show that conventional denoising
  filters are not consistently reliable as adversarial defenses for AMC:
  while they offer partial recovery against CW and EAD, several of them
  reduce accuracy below the undefended baseline on FGSM.

\item \textbf{Low-latency accuracy recovery across multiple attack families.}
  Adaptive-$K$ achieves the best recovery on four of five evaluated attacks,
  improving CW accuracy from 35.4\% to 71.3\% (+35.9\,pp) and EAD-L1 accuracy
  from 6.6\% to 70.8\% (+64.2\,pp), while degrading clean accuracy by less
  than 0.5\,pp and adding less than 0.1\,ms of latency per batch of 64
  samples in our GPU implementation.
\end{enumerate}

In our adaptive-$K$ method, the selected $K$ should be interpreted as the
number of \emph{dominant spectral bins} retained by the Top-$K$ operator,
not as a direct physical-bandwidth estimate. Different modulations may
therefore exhibit different $K$ values depending on the spectral concentration
of their energy under the recorded conditions, even when their nominal
occupied bandwidths are similar.

The remainder of this thesis is organized as follows. Chapter~\ref{ch:relatedwork}
reviews related work on adversarial AMC attacks, defense methods, and classical
signal-processing filters. Chapter~\ref{ch:architecture} presents the system
model, attack formulations, and threat model. Chapter~\ref{ch:method}
describes the proposed Adaptive-$K$ defense and the classical filter
baselines. Chapter~\ref{ch:evaluation} presents experimental results and
analysis. Chapter~\ref{ch:conclusion} concludes the thesis.
